% Basic stuff
\documentclass[a4paper,11pt]{article}
\usepackage[english]{babel}

% 3 column landscape layout with fewer margins
\usepackage[landscape, left=0.75cm, top=1cm, right=0.75cm, bottom=1.5cm, footskip=15pt]{geometry}
\usepackage{flowfram}
\ffvadjustfalse
\setlength{\columnsep}{1cm}
\Ncolumn{3}

% Reduce vertical spacing
\usepackage{titlesec}
\titlespacing*{\section}{0pt}{1ex plus .5ex minus .2ex}{0.5ex plus .2ex}
\titlespacing*{\subsection}{0pt}{0.75ex plus .5ex minus .2ex}{0.25ex plus .2ex}

\makeatletter
\g@addto@macro\normalsize{%
  \setlength\abovedisplayskip{3pt}%
  \setlength\belowdisplayskip{3pt}%
  \setlength\abovedisplayshortskip{2pt}%
  \setlength\belowdisplayshortskip{2pt}%
}
\makeatother

% Make enumerations more compact and adjust indentation
\usepackage{enumitem}
\setlist[itemize,1]{itemsep=1pt, topsep=2pt, left=0pt}
\setlist[itemize,2]{itemsep=1pt, topsep=1pt, left=10pt}
\setlist[enumerate,1]{itemsep=1pt, topsep=2pt, left=0pt}
\setlist[enumerate,2]{itemsep=1pt, topsep=1pt, left=10pt}

% define nice looking boxes
\usepackage[many]{tcolorbox}

% a base set, that is then customised
\tcbset {
  base/.style={
    boxrule=0mm,
    leftrule=1mm,
    left=1.75mm,
    arc=0mm, 
    fonttitle=\bfseries, 
    colbacktitle=black!10!white, 
    coltitle=black, 
    toptitle=0.75mm, 
    bottomtitle=0.25mm,
    title={#1},
    breakable,
    top=1pt,
    bottom=1pt
  }
}

\definecolor{iml}{rgb}{0.19, 0.62, 0.75}
\newtcolorbox{mainbox}[1]{
  colframe=iml, 
  base={#1}
}

\newtcolorbox{subbox}[1]{
  colframe=black!20!white,
  base={#1}
}

% Mathematical typesetting & symbols
\usepackage{amsthm, mathtools, amssymb} 
\usepackage{marvosym, wasysym}
\usepackage{derivative}
\allowdisplaybreaks

%Algorithms
\usepackage{algpseudocode}
\usepackage{algorithm}
% Tables
\usepackage{tabularx, multirow}
\usepackage{booktabs}
\usepackage{float}

% Make enumerations more compact
\usepackage{enumitem}
\setitemize{itemsep=0.5pt}
\setenumerate{itemsep=0.75pt}

% To include sketches & PDFs
\usepackage{graphicx}

% For hyperlinks
\usepackage{hyperref}
\hypersetup{
  colorlinks=true
}

% Metadata
\title{Cheatsheet Introduction to Machine Learning}
\author{Florian Affolter}
\date{Sommersession 2025}

% Math helper stuff
\def\R{\mathbb{R}}
\def\P{\mathbb{P}}
\def\E{\mathbb{E}}
\def\T{\mathsf{T}}

\renewcommand\thesection{\Roman{section}}
\usepackage{remreset}
\makeatletter
\@removefromreset{subsection}{section}
\renewcommand\thesubsection{\arabic{subsection}}
\makeatother

\begin{document}

% \maketitle

\section*{Preliminaries}
\subsection*{Linear Algebra}

A symmetric matrix with real entries is \emph{positive (semi-)definite} iff all its eigenvalues are positive (nonnegative).

\( A A^\T \) and \( A^\T A \) are symmetric and positive semi-definite. If \( m \leq n \) and \( \operatorname{rank}(A) = m \), then \( A A^\T \) is also positive definite.


\subsection*{Analysis}

\[
        D f(\mathbf{x_0}) = J_f(\mathbf{x_0}) \\= \begin{bsmallmatrix}
        \frac{\partial f_1}{\partial x_1}(\mathbf{x_0}) & \cdots & \frac{\partial f_1}{\partial x_n}(\mathbf{x_0}) \\
        \vdots & \ddots & \vdots \\
        \frac{\partial f_m}{\partial x_1}(\mathbf{x_0}) & \cdots & \frac{\partial f_m}{\partial x_n}(\mathbf{x_0})
    \end{bsmallmatrix}
\]

\[
    D f(\mathbf{x_0})[\mathbf{h}] = \langle \nabla f(\mathbf{x_0}), \mathbf{h} \rangle
\]

\[
    D (f \circ g)(\mathbf{x_0}) = (D f)(g(\mathbf{x_0})) D g(\mathbf{x_0})
\]

\[
    D^2 f(\mathbf{x_0}) = H_f(\mathbf{x_0}) \\= \begin{bsmallmatrix}
        \frac{\partial^2 f}{\partial x_1^2}(\mathbf{x_0}) & \cdots & \frac{\partial^2 f}{\partial x_n \partial x_1}(\mathbf{x_0}) \\
        \vdots & \ddots & \vdots \\
        \frac{\partial^2 f}{\partial x_1 \partial x_n}(\mathbf{x_0}) & \cdots & \frac{\partial^2 f}{\partial x_n^2}(\mathbf{x_0})
    \end{bsmallmatrix}
\]


\subsection*{Probability Theory}

\begin{mainbox}{CDF}
    \[
        F_X(x) = \P_X((-\infty, x]) = \P(X \leq x)
    \]
\end{mainbox}
Properties: \( \lim_{x \to \infty} F_X(x) = 1 \) and \( \lim_{x \to -\infty} F_X(x) = 0 \). It is monotonically increasing. It is right-continuous. \( \P(a < X \leq b) = F_X(b) - F_X(a) \).

\begin{mainbox}{PMF}
    \[
        p_X(x) = \P(X = x) = \P_X(\{ x \})
    \]
\end{mainbox}

\begin{mainbox}{PDF}
    If \( \exists \; p_X: \: \R \to [0, \infty) \), such that \[
        \P_X(I) = \P(X \in I) = \int_I p_X(x) \, \mathrm{d}x
    \] for all intervals \( I \) in \( \R \), we call it \emph{PDF} of \( X \).
\end{mainbox}
\[
    F_X(x) = \P(X \leq x) = \int_{-\infty}^x p_X(t) \mathrm{d}t
\]


\section*{Supervised Learning}

\subsection*{Linear Regression}

\begin{subbox}{Normal Equation}
    \[
        X^\T X \hat{w} = X^\T y
    \]
\end{subbox}

Unique solution for the normal equation (given that \( d \leq n \) and \( X \) is full-rank):
\[
    \hat{\omega} = \left( X^\T X \right)^{-1} X^\T y
\]

When over-parametrized, use pseudo-inverse:
\[
    \hat{w} = \left( X^\T X \right)^\dagger X^\T y
\]

If \( X \) is full-rank and \( n < d \), then \( y = Xw \) has infinitely many solutions. The minimum-norm solution is:
\[
    \hat{w} = X^\T \left( X X^\T \right)^{-1} y = \left( X^\T X \right)^\dagger X^\T y
\]


\subsection*{Optimization}

\begin{mainbox}{Convexity}
    A set \( \mathcal{W} \subseteq \R^d \) is convex if it contains the line segment between any two of its points
\end{mainbox}
\begin{enumerate}
    \item \( \alpha f + \beta g \) is convex if \( f, \, g \) are convex and \( \alpha, \, \beta \geq 0 \).
    \item \( f \circ g \) is convex if \( f \) is convex and \( g \) is affine, or if \( f \) is convex and non-decreasing and \( g \) is convex.
    \item The point-wise maximum of two convex functions is convex.
\end{enumerate}

If \( f \) is \emph{strongly convex}, at every point there is always a quadratic function that fits between \( f \) and the tangent of \( f \).
The notion of strong convexity excludes functions with linear parts.

\begin{subbox}{}
    Suppose \( f \) is differentiable and strongly convex. Then there exists unique minimizer \( \mathbf{w^*} \) of \( f \) and gradient desent with sufficiently small step size and arbitrary initialization \( \mathbf{w^0} \) satisfies \[
        \lim_{t \to \infty} \mathbf{w^t} = \mathbf{w^*} \text{.}
    \]
\end{subbox}


\subsection*{Model evaluation and selection}

\begin{subbox}{Error metrics}
    \begin{itemize}
        \item Model training
        \begin{itemize}
            \item \textbf{Training error (empirical risk):} \[
                \frac{1}{\lvert \mathcal{D} \rvert} \sum_{(x, y) \in \mathcal{D}} \ell(f(x), y)
            \] Used to train a model \( \hat{f}_\mathcal{D} \).
        \end{itemize}
    \end{itemize}
    \begin{itemize}
        \item Model evaluation
        \begin{itemize}
            \item \textbf{Expected estimation error:} \[
                \E_X \ell \left( \hat{f}_\mathcal{D}(X), f^*(X) \right)
            \] Tells us how good \( \hat{f}_D \) is. No access to \( f^* \)!
            \item \textbf{Generalization error (population risk):} \[
                \E_{X, Y} \ell \left( \hat{f}_\mathcal{D}(X), Y \right)
            \] Proxy for estimation error. No access to \( \P_{X, Y} \)!
            \item \textbf{Test error:} \[
                \frac{1}{\lvert \mathcal{D}_{\text{test}} \rvert} \sum_{(x, y) \in \mathcal{D}_\text{test}} \ell(f(x), y)
            \] Estimator of generalization error. Can be computed!
        \end{itemize}
    \end{itemize}
\end{subbox}

Even ground-truth function has strictly positive generalization error due to noise.

Under model \( y = f^*(x) + \epsilon \), and under square loss:
\[
    L \left( \hat{f}; \P_{X, Y} \right) = \E_X \left( \hat{f}(X) - f^*(X) \right)^2 + \sigma^2
\]


\subsection*{Bias-Variance Tradeoff and Regularization}

\begin{subbox}{Model Bias}
    \[
        \text{Bias}_\mathcal{D}^2 \left( \hat{f}_\mathcal{D}, \mathbf{x} \right) := \left( \E_\mathcal{D} \left[ \hat{f}_\mathcal{D}(\mathbf{x)} \right] - f^*(\mathbf{x)} \right)^2
    \]
    \[
        \text{Bias}_\mathcal{D}^2 \left( \hat{f}_\mathcal{D} \right) := \E_X \left[ \text{Bias}_\mathcal{D}^2 \left( \hat{f}_\mathcal{D}, X \right) \right]
    \]
\end{subbox}

\begin{subbox}{Model Variance}
    \begin{multline*}
        \text{Var}_\mathcal{D} \left( \hat{f}_\mathcal{D}(\mathbf{x}) \right) :=\\ \E_\mathcal{D} \left[ \left( \hat{f}_\mathcal{D}(\mathbf{x}) - \E_\mathcal{D} \left[ \hat{f}_\mathcal{D}(\mathbf{x}) \right] \right)^2 \right]
    \end{multline*}
    \[
        \text{Var}_\mathcal{D} \left( \hat{f}_\mathcal{D} \right) := \E_X \left[ \text{Var}_\mathcal{D} \left( \hat{f}_\mathcal{D}(X) \right) \right]
    \]
\end{subbox}

\begin{multline*}
    \E_\mathcal{D} \left[ L \left( \hat{f}_\mathcal{D}; \P_{X, Y} \right) \right] =\\ \text{Var}_\mathcal{D} \left( \hat{f}_\mathcal{D} \right) + \text{Bias}_\mathcal{D}^2 \left( \hat{f}_\mathcal{D} \right) + \sigma^2
\end{multline*}

As \( \lambda \) increases, the bias also increases, whereas the variance drops. The noise term stays constant.


\subsection*{Classification}

\begin{mainbox}{Support Vector Machine}
    \textbf{Hard-margin SVM:} \[
        \min_{w \in \R^d} \lVert \mathbf{w} \rVert_2
    \] such that \[
        y_i \langle \mathbf{w}, \mathbf{x_i} \rangle \geq 1, \; \forall i = 1, \dots, n
    \]

    \textbf{Soft-margin SVM:} \[
        \min_{w \in \R^d, \, \zeta \in \R^n} \left( \lVert \mathbf{w} \rVert_2^2 + \lambda \sum_{i = 1}^n \zeta_i \right)
    \] such that \[
        y_i \langle \mathbf{w}, \mathbf{x_i} \rangle \geq 1 - \zeta_i, \; \zeta_i \geq 0, \; \forall i = 1, \dots, n
    \]
\end{mainbox}

\[
    \text{FNR} := \frac{\text{\#FN}}{\text{\#}\{y = +1\}}, \ \text{FPR} := \frac{\text{\#FP}}{\text{\#}\{y = -1\}}
\]
\[
    \text{TPR}  := \frac{\text{\#TP}}{\text{\#}\{y = +1\}}, \ \text{TNR} := \frac{\text{\#TN}}{\text{\#}\{y = -1\}}
\]
TPR = \emph{power} = \emph{recall}.


\subsection*{Kernels}

\begin{mainbox}{Kernel}
    \( k: \: \R^d \times \R^d \to \R^d \) if it is symmetric and positive semi-definite.
\end{mainbox}
Feature maps can be \emph{composed}. Kernels can be \emph{added} in two ways: \[ k((\mathbf{x}, \mathbf{y}), (\mathbf{x'}, \mathbf{y'})) = k_1(\mathbf{x}, \mathbf{x'}) + k_2(\mathbf{y}, \mathbf{y'}) \text{,} \] \[ k(\mathbf{x}, \mathbf{x'}) = k_1(\mathbf{x}, \mathbf{x'}) + k_2(\mathbf{x}, \mathbf{x'}) \text{.} \] Kernels can be \emph{multiplied} in two ways: \[ k((\mathbf{x}, \mathbf{y}), (\mathbf{x'}, \mathbf{y'})) = k_1(\mathbf{x}, \mathbf{x'}) k_2(\mathbf{y}, \mathbf{y'}) \text{,} \] \[ k(\mathbf{x}, \mathbf{x'}) = k_1(\mathbf{x}, \mathbf{x'}) k_2(\mathbf{x}, \mathbf{x'}) \text{.} \]


\subsection*{Neural Networks}

\begin{mainbox}{Gradient}
    \[
        \nabla_\mathbf{x} f = \frac{\partial f}{\partial \mathbf{x}} = \begin{bsmallmatrix}
            \frac{\partial f}{\partial x_1}, & \dots, & \frac{\partial f}{\partial x_n}
        \end{bsmallmatrix}^\T
    \]
    \[
        \nabla_{\mathbf{W^{(l)}}} f = \frac{\partial f}{\partial \mathbf{W^{(l)}}} = \begin{bsmallmatrix}
            \frac{\partial f}{\partial w_{1, 1}^{(l)}} & \cdots & \frac{\partial f}{\partial w_{1, n}^{(l)}} \\
            \vdots & \ddots & \vdots \\
            \frac{\partial f}{\partial w_{m, 1}^{(l)}} & \cdots & \frac{\partial f}{\partial w_{m, n}^{(l)}}
        \end{bsmallmatrix}
    \]
\end{mainbox}


% \section{Unsupervised Learning}

% \subsection{Clustering}

% \subsection{Principal Component Analysis}


\section*{Probabilistic Methods}

\subsection*{Probabilistic Modeling \& Inference}

\begin{mainbox}{Bayes' Rule}
    \[
        \underbrace{p(\theta \; | \; \mathcal{D})}_\text{posterior belief} = \underbrace{\frac{p(\mathcal{D \; | \; \theta})}{p(\mathcal{D})}}_\text{update} \underbrace{p(\theta)}_\text{prior belief}
    \]
\end{mainbox}

\begin{mainbox}{MLE}
    \begin{align*}
        \hat{\theta}_\text{MLE} &:= \mathop{\arg \max}_{\theta \in \Theta} p(\mathcal{D} \; | \; \theta) \\
        &\overset{\text{i.i.d.}}{=} \mathop{\arg \max}_{\theta \in \Theta} \prod_{i = 1}^n p(\mathbf{x_i, y_i; \theta)} \\
        &= \mathop{\arg \min}_{\theta \in \Theta} \sum_{i = 1}^n -\log p(\mathbf{x_i}, y_i; \theta)
    \end{align*}
\end{mainbox}

\begin{mainbox}{MAP}
    \begin{align*}
        \hat{\theta}_\text{MAP} &:= \mathop{\arg \max}_{\theta \in \Theta} p(\theta \; | \; \mathcal{D}) \\
        &= \mathop{\arg \max}_{\theta \in \Theta} p(\mathcal{D} \; | \; \theta) p(\theta) \\
        &\overset{\text{i.i.d.}}{=} \mathop{\arg \max}_{\theta \in \Theta} \left( \prod_{i = 1}^n p(\mathbf{x_i, y_i; \theta)} \right) \cdot p(\theta) \\
        &= \mathop{\arg \min}_{\theta \in \Theta} \sum_{i = 1}^n -\log p(\mathbf{x_i}, y_i \; | \; \theta) - \log p(\theta)
    \end{align*}
    For discriminative models, equivalently:
    \[
        \mathop{\arg \min}_{\theta \in \Theta} \sum_{i = 1}^n -\log p(y_i \; | \; \mathbf{x_i}, \theta) - \log p(\theta)
    \]
\end{mainbox}


% \subsection{Gaussian Mixture Models}

% \subsection{Large Language Models}




\end{document}
